% ABSTRACT
% Last updated: 2026-02-21

Decentralized Autonomous Organizations (DAOs) represent a paradigm shift in organizational governance, yet their design space remains poorly understood. We present a comprehensive multi-agent simulation framework for analyzing governance mechanisms in DAOs, enabling systematic study of how design parameters affect collective decision-making outcomes. Our framework models heterogeneous agent populations with varying participation strategies, token distributions, and delegation behaviors, supporting multiple voting mechanisms including token-weighted, quadratic, and conviction voting. We additionally evaluate 14 historically conditioned DAO twins on a separately hashed temporal holdout. Mean composite similarity is 0.473; calibrated simulations outperform historical persistence for 5 of 14 DAOs and an uncalibrated simulator for 8 of 14, demonstrating heterogeneous generative fidelity rather than universal predictive accuracy.

Through Monte Carlo simulations (\totalruns{} runs across \experimentcount{} experiment configurations), we investigate fundamental questions in computational social choice: How do quorum thresholds affect proposal passage rates? What mitigation strategies effectively reduce governance capture? How do pipeline configurations impact proposal throughput? Our key findings include: (1) quadratic voting significantly increases decision margins ($d = -2.00$, $p < 0.0001$) while maintaining comparable capture risk; (2) quadratic thresholds of 250 tokens reduce whale influence by 43\% (from 0.449 to 0.256) without sacrificing governance throughput; (3) temp-check stages are critical for pipeline health, filtering low-quality proposals while passing viable initiatives; (4) treasury stabilization halves volatility with a minor growth cost; (5) inter-DAO cooperation succeeds at 21--23\% but requires explicit fairness mechanisms for asymmetric partnerships.

We release our simulation framework as open-source software to enable reproducible governance research and evidence-based mechanism design.
